% Generated draft: requires mathematical and source review.
\documentclass[11pt]{article}
\usepackage[margin=25mm]{geometry}
\usepackage{fontspec}
\setmainfont{Latin Modern Roman}
\usepackage{amsmath,amssymb,mathtools,unicode-math}
\setmathfont{Latin Modern Math}
\usepackage{enumitem,needspace,xurl,hyperref}
\setlength{\parindent}{0pt}
\setlength{\parskip}{5pt}
\setlength{\emergencystretch}{3em}
\newcommand{\N}{\mathbb{N}}
\newcommand{\Z}{\mathbb{Z}}
\newcommand{\Q}{\mathbb{Q}}
\newcommand{\R}{\mathbb{R}}
\newcommand{\C}{\mathbb{C}}
\newcommand{\F}{\mathbb{F}}
\newcommand{\A}{\mathbb{A}}
\newcommand{\PP}{\mathbb{P}}
\newcommand{\E}{\mathbb{E}}
\newcommand{\eps}{\varepsilon}
\newcommand{\dd}{\,\mathrm{d}}
\begin{document}
\section*{UnsolvedMath ID: 20003142}
\textbf{Source problem code:} catalogue entry
\subsection{Definitions and mathematical statement}
The source asks for an approach to homotopy groups adapted to Alexandrov spaces, without specifying desired axioms or quantitative features. The supplied partial answer uses Lipschitz representatives and a quantitative refinement of ordinary homotopy groups.

Let \((X,x_0)\) be a pointed finite-dimensional Alexandrov space, in the usual length-space, lower-curvature-bound convention, and let \(k\geq1\). Equip \(S^k\) with the unit round geodesic metric and a basepoint \(*\). Define \(\pi_k^{\mathrm{Lip}}(X,x_0)\) using based Lipschitz maps \(S^k\to X\), modulo based Lipschitz homotopies with the product metric. The aid reports published comparison isomorphisms
\[
 \pi_k^{\mathrm{Lip}}(X,x_0)\cong\pi_k(X,x_0).
\]
The precise comparison theorem, including completeness and boundary conventions, is not cited in full.

Via this comparison, for \(\alpha\in\pi_k(X,x_0)\) define
\[
 \sigma_{k,X}(\alpha)=\inf\{\operatorname{Lip}(f):f:(S^k,*)\to(X,x_0)
 \text{ Lipschitz and }[f]=\alpha\}.
\]
The aid states that this is finite for every class and zero for the trivial class. For every based Lipschitz map \(h:(X,x_0)\to(Y,y_0)\) between spaces in this setting,
\[
 \sigma_{k,Y}(h_*\alpha)\leq\operatorname{Lip}(h)\,\sigma_{k,X}(\alpha).
\]
The proposed size is not asserted a group norm with an exact triangle inequality. The reported product control is, with a suitable dimension-dependent Lipschitz pinch-map constant \(C_k\),
\[
 \sigma_{k,X}(\alpha\beta)\leq C_k\bigl(\sigma_{k,X}(\alpha)+\sigma_{k,X}(\beta)\bigr).
\]
Here multiplication is the usual possibly nonabelian product for \(k=1\), and may be written additively for \(k\geq2\). The exact chosen pinch map and constant are not supplied.

For the positive lower bound, assume \(U=B(x_0,\rho)\), \(\rho>0\), admits a strong local Lipschitz contraction to \(x_0\), fixing \(x_0\) and preserving every smaller concentric ball. In particular, the intended homotopy \(H\) has \(H(z,0)=z\), \(H(z,1)=x_0\), \(H(x_0,t)=x_0\), and \(H(B(x_0,r)\times[0,1])\subseteq B(x_0,r)\) for \(r<\rho\), with the Lipschitz regularity needed on compact restrictions. Since the based unit sphere has radius at most \(\pi\) from \(*\), the aid's bound is
\[
 \alpha\ne0\quad\Longrightarrow\quad\sigma_{k,X}(\alpha)\geq\rho/\pi.
\]
The normalization of the domain metric is essential to that constant.

The associated fixed-degree homotopy systole is
\[
 \operatorname{sys}_k(X,x_0)=\inf_{\alpha\ne0}\sigma_{k,X}(\alpha),
\]
with value \(+\infty\) when \(\pi_k=0\). Under the stated contraction it is at least \(\rho/\pi\), and is finite if there is a nonzero class. For every metric rescaling \(d_X\mapsto c\,d_X\), \(c>0\), both \(\sigma_{k,X}\) and this systole scale by \(c\). The same uniform lower bound applies if a global systole is defined by taking the infimum across all nontrivial degrees. These are the supplied quantitative conclusions, conditional on the contraction and comparison hypotheses, rather than a new characterization of Alexandrov homotopy types.
\subsection{Short English statement}
The supplied approach transfers Lipschitz representative sizes to ordinary homotopy classes. It gives functorial bounds, controlled products, and a positive homotopy systole under a specified strong local contraction. The constant \(\rho/\pi\) uses the unit round sphere and is not independent of domain normalization.
\subsection{Sources}
\begin{itemize}
\item Local UnsolvedMath registry, permanent ID \(20003142\), exact source code \texttt{catalogue entry}; canonical \texttt{aim-topology-notes.json}, \texttt{notes[229]}, source tag \texttt{problem}.
\item AIM workshop list ``Open Problems in Non-negative Sectional Curvature,'' workshop ``Manifolds with nonnegative sectional curvature,'' Problem 17: \url{https://aimath.org/WWN/nnsectcurvature/nnsectcurvature.pdf}. The supplied verification calls the canonical question exact. The PDF was not fetched.
\item The research summary refers to published work comparing Lipschitz and ordinary homotopy groups of finite-dimensional Alexandrov spaces, but supplies no authors, titles, identifiers, or full authoritative citations. The original literature field is empty.
\item Research attempt 1: result \texttt{partial\_result}, problem status at run \texttt{partially\_solved}, rigor label \texttt{proved}, novelty confidence \texttt{low}; classification \texttt{PARTIAL-PROGRESS}, registry status \texttt{partially\_solved}. Created and updated 14 August 2026. These are supplied assessments requiring review.
\item Attribution: Ulam AI / UnsolvedMath Contributors, CC BY 4.0. The machine-generated research aid is not peer reviewed and requires independent verification before citation.
\end{itemize}
\end{document}
